In December to February the three super El Niño winters were milder than neutral winters in Finland, Sweden, Norway and Denmark once the warming trend is removed (Nordic mean \nordicTSvNdiff\,\degC, 95\,\% interval \nordicTSvNlo{} to \nordicTSvNhi; Finland \fiTSvNdiff\,\degC, \fiTSvNlo{} to \fiTSvNhi), with more thaw days and slightly less snow, and Iceland was the exception with a small cooling. Against ordinary El Niño winters the temperature difference was smaller (Nordic \nordicTSvEdiff\,\degC, \nordicTSvElo{} to \nordicTSvEhi). Finland was the only country with a clear precipitation surplus (\fiPSvNdiff\,mm over the season, \fiPSvNlo{} to \fiPSvNhi). None of these differences reached $p < 0.05$ in a permutation test (Nordic temperature $p = \nordicTSvNp$; Finnish precipitation $p = \fiPSvNp$), the January to March window showed no signal at all, and the ONI had no monotonic relation with any metric (Spearman $|\rho| \le 0.15$). In absolute terms the three super winters were ordinary Finnish winters, not mild ones; the mildness appears only relative to the trend of their own era.
